% !TeX encoding = UTF-8
\documentclass[institutional,screen,none]{ua-engineering-v6-beamer}
% !TeX encoding = UTF-8
\UASetUniversity{Universidad Austral}
\UASetFaculty{Facultad de Ingeniería}
\UASetDepartment{Departamento de Computación}
\UASetCampus{Pilar}
\UASetCareer{Ingeniería Informática}
\UASetCommission{Comisión A}
\UASetProfessor{Equipo docente}
\UASetSemester{Segundo cuatrimestre 2026}
\UASetCourse{Sistemas Distribuidos}
\UASetDate{2026}

\UASetSemester{Segundo cuatrimestre 2026}
\UASetDate{2026}

\tikzset{
  uaBox/.style={draw=UAIngNavy,rounded corners=1.5mm,thick,fill=white,
    minimum height=0.75cm,align=center,font=\scriptsize,inner sep=4pt},
  uaBoxO/.style={uaBox,draw=UAIngOrange,fill=UAIngPaper},
  uaSoft/.style={uaBox,draw=UAIngRule,fill=UAIngPaper},
  uaArrow/.style={->,very thick,>=Latex,draw=UAIngNavy},
  uaArrowO/.style={->,very thick,>=Latex,draw=UAIngOrange}
}


% Etiqueta pedagógica visible. Se define localmente para no modificar el .sty
% compartido ni el build de GitHub.
\newcommand{\UASlideTag}[2]{%
  \begin{tikzpicture}[remember picture,overlay]
    \node[anchor=north east,rounded corners=1mm,inner xsep=3pt,inner ysep=1.5pt,
      draw=#2,fill=#2!8,text=#2,font=\fontsize{5.8}{6.4}\selectfont\bfseries]
      at ([xshift=-0.72cm,yshift=-1.08cm]current page.north east) {#1};
  \end{tikzpicture}%
}

\UASetClassNumber{09}
\UASetClassTitle{Arquitecturas distribuidas reales}
\title{\UACourse}
\subtitle{Clase 09 --- Monolito modular, microservicios, eventos, outbox y sagas}
\author{\UAProfessor}
\date{\UADate}

\begin{document}
{
\setbeamertemplate{footline}{}
\begin{frame}[plain]
\titlepage
\end{frame}
}

\begin{frame}{Cómo vamos a trabajar hoy}
\UASlideTag{NÚCLEO}{UAIngNavy}
\begin{center}
\begin{tikzpicture}[node distance=0.45cm]
\node[uaBoxO,text width=2.7cm] (a) {Intuición\\edificios y caminos};
\node[uaBox,text width=2.7cm,right=of a] (b) {Fronteras\\dominio y ownership};
\node[uaBoxO,text width=2.7cm,right=of b] (c) {Comunicación\\RPC y eventos};
\node[uaBox,text width=2.7cm,right=of c] (d) {Consistencia\\outbox y sagas};
\draw[uaArrow] (a)--(b); \draw[uaArrowO] (b)--(c); \draw[uaArrow] (c)--(d);
\end{tikzpicture}
\end{center}
\begin{block}{Caso conductor}
Asteria Online debe separar cuentas, mundo, combate, inventario, pagos, ranking y chat. La pregunta no es cuántos servicios crear, sino qué frontera compra autonomía suficiente para justificar el costo distribuido.
\end{block}
\end{frame}

\begin{frame}{Evidencias de aprendizaje que produciremos}
\UASlideTag{NÚCLEO}{UAIngNavy}
\begin{columns}[T]
\column{0.48\textwidth}
\begin{block}{Durante la clase}
\begin{itemize}
\item evaluar una frontera con capacidad, invariante y owner;
\item calcular la prima de una cadena sincrónica;
\item mover un crash por el dual write.
\end{itemize}
\end{block}
\column{0.48\textwidth}
\begin{block}{Al cerrar}
\begin{itemize}
\item diseñar outbox e idempotencia;
\item modelar una saga durable y sus compensaciones;
\item proponer una migración reversible con métricas.
\end{itemize}
\end{block}
\end{columns}
\begin{uakeybox}
Una frontera solo se aprueba si compra autonomía medible y si el equipo puede operar las fallas que introduce.
\end{uakeybox}
\end{frame}


\UASectionSlide{De la consistencia del dato\\a la frontera del sistema}

\begin{frame}{Puente desde la clase 8}
\UASlideTag{NÚCLEO}{UAIngNavy}
\begin{itemize}
\item Cada dato exige una decisión distinta de consistencia, disponibilidad y latencia.
\item Una arquitectura distribuye esas decisiones entre módulos, servicios y equipos.
\item Mover una frontera cambia qué invariantes son locales y cuáles requieren protocolo.
\end{itemize}
\begin{uakeybox}
La pregunta de hoy es: ``¿qué complejidad distribuida estamos comprando y qué autonomía real recibimos a cambio?''.
\end{uakeybox}
\end{frame}

\begin{frame}{Diagnóstico de entrada: vote antes de discutir}
\UASlideTag{ACTIVIDAD}{UAIngOrange}
\small
\begin{enumerate}
\item Un monolito es necesariamente una bola de barro.
\item Un microservicio es una clase de negocio expuesta por HTTP.
\item Publicar un evento después de guardar en la base resuelve el dual write.
\item Una compensación de saga equivale a un rollback perfecto.
\item Más servicios implican automáticamente más disponibilidad.
\end{enumerate}
\begin{block}{Consigna}
Vote V/F. Volveremos al final con casos, timelines y métricas.
\end{block}
\end{frame}

\begin{frame}{Caso conductor: Asteria Online}
\UASlideTag{NÚCLEO}{UAIngNavy}
\begin{center}
\begin{tikzpicture}[node distance=0.30cm]
\node[uaBoxO,text width=1.75cm] (acc) {Cuentas};
\node[uaBox,text width=1.65cm,right=of acc] (world) {Mundo};
\node[uaBoxO,text width=1.65cm,right=of world] (combat) {Combate};
\node[uaBox,text width=1.65cm,right=of combat] (inv) {Inventario};
\node[uaBoxO,text width=1.75cm,right=of inv] (pay) {Pagos};
\node[uaBox,text width=1.65cm,right=of pay] (rank) {Ranking};
\end{tikzpicture}
\end{center}
\begin{itemize}
\item Combate necesita baja latencia y autoridad por zona.
\item Inventario debe impedir duplicación de objetos.
\item Ranking tolera atraso.
\item Pagos requieren trazabilidad e idempotencia.
\item Chat puede degradarse sin detener el mundo.
\end{itemize}
\end{frame}

\begin{frame}{Predicción individual: ¿dónde pondría la primera frontera?}
\UASlideTag{ACTIVIDAD}{UAIngOrange}
\begin{uainfo}
Elija una capacidad de Asteria que mantendría dentro del monolito y una que separaría. Escriba una sola razón de negocio y una sola consecuencia operativa para cada decisión.
\end{uainfo}
\vspace{0.25cm}
Guarde la respuesta: volveremos a ella después de estudiar ownership, prima distribuida, outbox y sagas.
\end{frame}


\begin{frame}{Cinco preguntas antes de dibujar cajas}
\UASlideTag{NÚCLEO}{UAIngNavy}
\begin{center}
\begin{tikzpicture}[node distance=0.18cm]
\node[uaBoxO,text width=1.98cm] (q1) {¿Qué capacidad\\de negocio?};
\node[uaBox,text width=1.98cm,right=of q1] (q2) {¿Quién posee\\el dato?};
\node[uaBoxO,text width=1.98cm,right=of q2] (q3) {¿Qué invariante\\debe ser local?};
\node[uaBox,text width=1.98cm,right=of q3] (q4) {¿Qué falla nueva\\aparece?};
\node[uaBoxO,text width=1.98cm,right=of q4] (q5) {¿Qué autonomía\\se gana?};
\draw[uaArrow] (q1)--(q2); \draw[uaArrowO] (q2)--(q3); \draw[uaArrow] (q3)--(q4); \draw[uaArrowO] (q4)--(q5);
\end{tikzpicture}
\end{center}
\end{frame}

\UASectionSlide{Primero la intuición: edificios,\\caminos y propietarios}

\begin{frame}{Monolito modular: un edificio con habitaciones claras}
\UASlideTag{NÚCLEO}{UAIngNavy}
\begin{center}
\begin{tikzpicture}[x=1cm,y=1cm]
\draw[thick,UAIngNavy,rounded corners] (0,0) rectangle (9,3.5);
\draw[thick,UAIngOrange] (3,0)--(3,3.5);
\draw[thick,UAIngNavy] (6,0)--(6,3.5);
\node[align=center] at (1.5,1.75){Cuentas\\módulo};
\node[align=center] at (4.5,1.75){Mundo\\módulo};
\node[align=center] at (7.5,1.75){Inventario\\módulo};
\node[above,font=\bfseries] at (4.5,3.6){una unidad de despliegue};
\end{tikzpicture}
\end{center}
\begin{block}{Intuición}
Compartir edificio simplifica transacciones, debugging y despliegue. La calidad depende de paredes internas reales, no del número de procesos.
\end{block}
\end{frame}

\begin{frame}{Microservicios: edificios separados conectados por caminos}
\UASlideTag{NÚCLEO}{UAIngNavy}
\begin{center}
\begin{tikzpicture}[node distance=0.8cm]
\node[uaBoxO,text width=2.2cm] (a) {Cuentas\\equipo A};
\node[uaBox,text width=2.2cm,right=1.4cm of a] (b) {Mundo\\equipo B};
\node[uaBoxO,text width=2.2cm,right=1.4cm of b] (c) {Inventario\\equipo C};
\draw[uaArrow] (a)--node[above,font=\scriptsize]{red}(b);
\draw[uaArrowO] (b)--node[above,font=\scriptsize]{red}(c);
\end{tikzpicture}
\end{center}
\begin{uakeybox}
Separar compra autonomía de despliegue, escala y ownership. El precio son caminos: latencia, fallas parciales, contratos, observabilidad y seguridad.
\end{uakeybox}
\end{frame}

\begin{frame}{Bounded context: cada barrio tiene su diccionario}
\UASlideTag{NÚCLEO}{UAIngNavy}
\begin{center}
\begin{tabular}{p{0.21\textwidth}p{0.31\textwidth}p{0.31\textwidth}}
\toprule
Palabra & Contexto & Significado \\
\midrule
Jugador & Combate & posición, vida, cooldown \\
Jugador & Facturación & cuenta, moneda, impuestos \\
Objeto & Inventario & propiedad, cantidad, versión \\
Objeto & Catálogo & nombre, rareza, precio base \\
\bottomrule
\end{tabular}
\end{center}
\begin{block}{Intuición}
Forzar un modelo universal es como imponer un único diccionario a toda la ciudad. Traducir entre contextos hace explícitas las diferencias.
\end{block}
\end{frame}

\begin{frame}{Un bounded context no obliga a crear un proceso}
\UASlideTag{NÚCLEO}{UAIngNavy}
\begin{itemize}
\item Es primero una frontera semántica: vocabulario, modelo e invariantes coherentes.
\item Puede implementarse como módulo dentro de un monolito o como servicio desplegable.
\item La distribución física se justifica después por escala, cadencia, ownership o aislamiento.
\end{itemize}
\begin{alertblock}{Error de diseño}
Convertir cada sustantivo o cada contexto en un microservicio antes de demostrar el beneficio produce fronteras costosas y frágiles.
\end{alertblock}
\end{frame}


\begin{frame}{Ownership: una oficina firma la verdad}
\UASlideTag{NÚCLEO}{UAIngNavy}
\begin{center}
\begin{tikzpicture}[node distance=0.65cm]
\node[uaBoxO,text width=2.5cm] (owner) {Inventario\\owner de items};
\node[uaBox,text width=2.2cm,right=1.2cm of owner] (combat) {Combate\\consulta/equipa};
\node[uaBox,text width=2.2cm,above right=0.4cm and 1.0cm of combat] (market) {Mercado\\reserva/transfiere};
\node[uaBox,text width=2.2cm,below right=0.4cm and 1.0cm of combat] (rank) {Ranking\\solo observa};
\draw[uaArrow] (combat)--(owner); \draw[uaArrowO] (market)--(owner); \draw[uaArrow] (owner)--(rank);
\end{tikzpicture}
\end{center}
\begin{block}{Regla intuitiva}
Un dato crítico necesita un owner claro. Los demás piden acciones o consumen hechos; no escriben silenciosamente las tablas internas del owner.
\end{block}
\end{frame}

\begin{frame}{RPC: una llamada telefónica que exige respuesta ahora}
\UASlideTag{NÚCLEO}{UAIngNavy}
\begin{center}
\begin{tikzpicture}[node distance=0.8cm]
\node[uaBoxO] (a) {Checkout};
\node[uaBox,text width=2.5cm,right=2.0cm of a] (b) {Pagos};
\draw[uaArrow] (a)--node[above,font=\scriptsize]{authorize}(b);
\draw[uaArrowO] (b)--node[below,font=\scriptsize]{sí / no / timeout}(a);
\end{tikzpicture}
\end{center}
\begin{itemize}
\item Causalidad explícita y respuesta en el camino crítico.
\item Latencia y disponibilidad de la dependencia se heredan.
\item Timeouts, retries e idempotencia vuelven a aparecer.
\end{itemize}
\end{frame}

\begin{frame}{Evento: publicar un hecho en un tablero durable}
\UASlideTag{NÚCLEO}{UAIngNavy}
\begin{center}
\begin{tikzpicture}[node distance=0.55cm]
\node[uaBoxO] (p) {Orders};
\node[uaBox,text width=2.2cm,right=of p] (log) {\texttt{OrderConfirmed}};
\node[uaBoxO,text width=2.0cm,above right=0.3cm and 0.8cm of log] (n) {Notificación};
\node[uaBox,text width=2.0cm,right=0.8cm of log] (a) {Analytics};
\node[uaBoxO,text width=2.0cm,below right=0.3cm and 0.8cm of log] (r) {Ranking};
\draw[uaArrow] (p)--(log); \draw[uaArrowO] (log)--(n); \draw[uaArrow] (log)--(a); \draw[uaArrowO] (log)--(r);
\end{tikzpicture}
\end{center}
\begin{block}{Intuición}
El productor declara un hecho sin esperar a todos los lectores. Gana desacoplamiento temporal; paga orden parcial, duplicados, evolución de contratos y debugging distribuido.
\end{block}
\end{frame}

\begin{frame}{Outbox: anotar la carta en el mismo libro antes de enviarla}
\UASlideTag{NÚCLEO}{UAIngNavy}
\begin{center}
\begin{tikzpicture}[node distance=0.6cm]
\node[uaBoxO,text width=2.2cm] (svc) {Servicio};
\node[uaBox,text width=2.4cm,right=of svc] (db) {Transacción local\\estado + outbox};
\node[uaBoxO,text width=2.1cm,right=of db] (courier) {Mensajero\\repetible};
\node[uaBox,text width=2.1cm,right=of courier] (broker) {Broker};
\draw[uaArrow] (svc)--(db); \draw[uaArrowO] (db)--(courier); \draw[uaArrow] (courier)--(broker);
\end{tikzpicture}
\end{center}
\begin{uakeybox}
El outbox no hace magia exactly-once: vuelve durable la intención de publicar y permite retry controlado con IDs y deduplicación.
\end{uakeybox}
\end{frame}

\begin{frame}{Saga: un itinerario con acciones compensatorias}
\UASlideTag{NÚCLEO}{UAIngNavy}
\begin{center}
\begin{tikzpicture}[node distance=0.45cm]
\node[uaBoxO,text width=2.1cm] (s) {Reservar\\stock};
\node[uaBox,text width=2.1cm,right=of s] (p) {Cobrar\\pago};
\node[uaBoxO,text width=2.1cm,right=of p] (e) {Emitir\\entrega};
\node[uaSoft,text width=2.1cm,below=0.8cm of s] (cs) {Liberar stock};
\node[uaSoft,text width=2.1cm,below=0.8cm of p] (cp) {Reembolsar};
\draw[uaArrow] (s)--(p); \draw[uaArrowO] (p)--(e);
\draw[uaArrow,dashed] (p)--(cs); \draw[uaArrowO,dashed] (e)--(cp);
\end{tikzpicture}
\end{center}
\begin{block}{Intuición}
Una compensación es una nueva acción de negocio. Puede llegar tarde, fallar o no borrar efectos físicos ya visibles.
\end{block}
\end{frame}

\UASectionSlide{Formalizar fronteras: autonomía,\\invariantes y prima distribuida}

\begin{frame}{Arquitectura como asignación de responsabilidades}
\UASlideTag{NÚCLEO}{UAIngNavy}
Una frontera está bien defendida si declara:
\begin{enumerate}
\item capacidad de negocio;
\item owner del dato y del on-call;
\item invariantes locales;
\item contrato de entrada y salida;
\item fallas toleradas;
\item autonomía ganada;
\item costo operacional aceptado.
\end{enumerate}
\begin{uawarn}
``Kafka + Kubernetes + microservicios'' es una lista de tecnologías, no una arquitectura.
\end{uawarn}
\end{frame}

\begin{frame}{Monolito modular no significa bola de barro}
\UASlideTag{NÚCLEO}{UAIngNavy}
\begin{center}
\begin{tabular}{p{0.27\textwidth}p{0.30\textwidth}p{0.30\textwidth}}
\toprule
Dimensión & Modular monolith & Big ball of mud \\
\midrule
Dependencias & reglas explícitas & cualquier módulo llama a todo \\
Ownership & módulos y equipos claros & responsabilidad difusa \\
Datos & acceso mediado & tablas compartidas sin límites \\
Cambios & contratos internos & efectos impredecibles \\
\bottomrule
\end{tabular}
\end{center}
\begin{block}{Caso real}
Shopify ha documentado su trabajo para modularizar un gran monolito Rails y hacer visibles dependencias y componentes sin asumir que microservicios sean la única solución.
\end{block}
\end{frame}

\begin{frame}{Cuándo el monolito modular es una ventaja}
\UASlideTag{NÚCLEO}{UAIngNavy}
\begin{itemize}
\item Equipo pequeño y dominio todavía cambiante.
\item Invariantes atraviesan varios módulos y necesitan transacción local.
\item Escalado relativamente homogéneo.
\item El costo de plataforma distribuida superaría el beneficio.
\item Se necesita aprender las fronteras antes de congelarlas en red.
\end{itemize}
\begin{uakeybox}
La simplicidad es una propiedad de arquitectura que debe defenderse, no una señal de inmadurez.
\end{uakeybox}
\end{frame}

\begin{frame}{Microservicio útil: autonomía verificable}
\UASlideTag{NÚCLEO}{UAIngNavy}
\begin{itemize}
\item Capacidad de negocio coherente.
\item Owner claro de comportamiento y datos.
\item Despliegue y rollback independientes en la mayoría de los cambios.
\item Escalado diferente y justificable.
\item Contratos explícitos y observabilidad propia.
\end{itemize}
\begin{alertblock}{No es}
Una clase convertida en proceso, una tabla envuelta en HTTP o un servicio genérico que todos llaman para todo.
\end{alertblock}
\end{frame}

\begin{frame}{La prima distribuida}
\UASlideTag{NÚCLEO}{UAIngNavy}
\begin{center}
\begin{tikzpicture}[node distance=0.45cm]
\node[uaBoxO,text width=2.2cm] (gain) {Autonomía\\escala / despliegue};
\node[uaBox,text width=1.98cm,right=of gain] (rpc) {RPC\\timeouts / retries};
\node[uaBoxO,text width=1.98cm,right=of rpc] (data) {Datos\\consistencia};
\node[uaBox,text width=1.98cm,right=of data] (ops) {Operación\\tracing / on-call};
\node[uaBoxO,text width=1.98cm,right=of ops] (sec) {Seguridad\\identidad / secretos};
\draw[uaArrow] (gain)--(rpc); \draw[uaArrowO] (rpc)--(data); \draw[uaArrow] (data)--(ops); \draw[uaArrowO] (ops)--(sec);
\end{tikzpicture}
\end{center}
\begin{block}{Decisión}
Se paga la prima solo cuando autonomía, escala o organización producen un beneficio real y medible.
\end{block}
\end{frame}

\begin{frame}{Disponibilidad compuesta: cada dependencia entra al camino crítico}
\UASlideTag{NÚCLEO}{UAIngNavy}
Si cuatro servicios independientes tienen disponibilidad $0.999$ y todos son necesarios:
\[
A_{camino}=0.999^4\approx 0.996
\]
\begin{itemize}
\item El cálculo es simplificado: las fallas pueden estar correlacionadas.
\item La latencia de cola también se compone.
\item Caching, degradación y asincronía pueden retirar dependencias del camino crítico.
\end{itemize}
\begin{uawarn}
Más servicios no implican automáticamente mayor disponibilidad; pueden crear más puntos de espera y falla.
\end{uawarn}
\end{frame}

\begin{frame}{Monolito distribuido: pagar el costo sin ganar autonomía}
\UASlideTag{NÚCLEO}{UAIngNavy}
\begin{center}
\begin{tikzpicture}[node distance=0.45cm]
\node[uaBoxO] (a) {A}; \node[uaBox,right=of a] (b) {B}; \node[uaBoxO,right=of b] (c) {C}; \node[uaBox,right=of c] (d) {D};
\draw[uaArrow] (a)--(b); \draw[uaArrowO] (b)--(c); \draw[uaArrow] (c)--(d);
\node[uaSoft,text width=8.2cm,below=0.9cm of $(b)!0.5!(c)$] {base compartida + despliegue coordinado + cadena sincrónica};
\end{tikzpicture}
\end{center}
\begin{itemize}
\item Servicios separados, pero cada cambio exige coordinar todos.
\item Una base compartida permite acoplamiento invisible.
\item Si falla un servicio, todo el flujo se detiene.
\end{itemize}
\end{frame}

\begin{frame}{Caso real: modularidad antes de extraer}
\UASlideTag{NÚCLEO}{UAIngNavy}
\small
Shopify ha descrito un enfoque de monolito modular con componentes y herramientas para hacer cumplir dependencias. También ha extraído componentes cuando una presión concreta —por ejemplo despliegue geográfico o performance— justificó la separación.
\begin{itemize}
\item Modificar fronteras dentro del proceso es más barato que hacerlo sobre red.
\item La extracción es una decisión puntual, no una migración religiosa.
\item Métrica: lead time, frecuencia de despliegue, rendimiento y blast radius.
\end{itemize}
{\scriptsize Fuentes: Shopify Engineering, artículos sobre modular monolith y componentización.}
\end{frame}

\begin{frame}{Caso real: Netflix y autonomía de escala}
\UASlideTag{RESPALDO}{UAIngMuted}
\small
Netflix ha documentado sistemas donde servicios y pipelines tienen cargas, ciclos de cambio y estrategias de escalado diferentes. El beneficio aparece junto con una plataforma madura de despliegue, observabilidad y ownership.
\begin{itemize}
\item La independencia de despliegue requiere contratos compatibles.
\item Escala técnica debe acompañar escala organizacional.
\item Microservicios sin plataforma trasladan complejidad al equipo de producto.
\end{itemize}
{\scriptsize Fuente: Netflix TechBlog, casos de pipelines y event-driven systems.}
\end{frame}

\UASectionSlide{Comunicación: llamada,\\hecho y contrato}

\begin{frame}{Comando, evento y consulta no son sinónimos}
\UASlideTag{NÚCLEO}{UAIngNavy}
\begin{center}
\begin{tabular}{p{0.20\textwidth}p{0.31\textwidth}p{0.33\textwidth}}
\toprule
Mensaje & Intención & Ejemplo Asteria \\
\midrule
Comando & pedir al owner que haga algo & \texttt{ReserveItem} \\
Evento & declarar un hecho ocurrido & \texttt{ItemReserved} \\
Consulta & solicitar información sin cambiar estado & \texttt{GetInventory} \\
\bottomrule
\end{tabular}
\end{center}
\begin{alertblock}{Error frecuente}
Nombrar un comando en pasado no lo convierte en evento. La semántica depende de quién puede rechazarlo y quién posee la decisión.
\end{alertblock}
\end{frame}

\begin{frame}{Elegir RPC o evento con una pregunta simple}
\UASlideTag{NÚCLEO}{UAIngNavy}
\begin{center}
\begin{tikzpicture}[node distance=0.7cm]
\node[uaBoxO,text width=2.5cm] (q) {¿Necesito una respuesta\\para continuar ahora?};
\node[uaBox,text width=2.6cm,below left=0.8cm and 1.1cm of q] (yes) {Sí\\RPC / comando};
\node[uaBoxO,text width=2.6cm,below right=0.8cm and 1.1cm of q] (no) {No\\evento / trabajo async};
\draw[uaArrow] (q)--node[left,font=\scriptsize]{camino crítico}(yes);
\draw[uaArrowO] (q)--node[right,font=\scriptsize]{desacoplar}(no);
\end{tikzpicture}
\end{center}
\begin{block}{Después preguntar}
¿Qué ocurre si el receptor tarda, no existe, procesa dos veces o cambia su contrato?
\end{block}
\end{frame}

\begin{frame}{Dual write: dos sistemas, una ventana de inconsistencia}
\UASlideTag{NÚCLEO}{UAIngNavy}
\begin{center}
\begin{tikzpicture}[x=1.1cm,y=0.9cm]
\draw[->,thick] (0,0)--(8.5,0) node[right]{tiempo};
\foreach \x/\lab in {1/{guardar orden},3/{COMMIT DB},4.5/{CRASH},6.5/{publicar evento}}{
  \draw (\x,-0.12)--(\x,0.12); \node[above,align=center,font=\scriptsize] at (\x,0.18){\lab};
}
\draw[very thick,UAIngOrange] (4.5,0.9)--(4.5,-0.6);
\node[uaSoft,text width=3.0cm] at (2.6,-1.0){estado confirmado};
\node[uaSoft,text width=3.0cm] at (6.5,-1.0){evento ausente};
\end{tikzpicture}
\end{center}
\begin{uawarn}
Publicar primero solo invierte el problema: consumidores ven un evento de una transacción que podría abortar.
\end{uawarn}
\end{frame}

\begin{frame}{Transactional outbox: frontera local coherente}
\UASlideTag{NÚCLEO}{UAIngNavy}
\begin{center}
\begin{tikzpicture}[node distance=0.6cm]
\node[uaBoxO,text width=2.2cm] (cmd) {\texttt{ConfirmOrder}};
\node[uaBox,text width=2.4cm,right=of cmd] (tx) {Una transacción\\order + outbox};
\node[uaBoxO,text width=2.1cm,right=of tx] (relay) {Relay\\at-least-once};
\node[uaBox,text width=2.2cm,right=of relay] (broker) {Broker\\event\_id};
\draw[uaArrow] (cmd)--(tx); \draw[uaArrowO] (tx)--(relay); \draw[uaArrow] (relay)--(broker);
\end{tikzpicture}
\end{center}
\begin{itemize}
\item Estado e intención de publicar confirman juntos.
\item El relay puede duplicar y debe reintentar.
\item Consumidores necesitan deduplicación y contrato versionado.
\end{itemize}
\end{frame}

\begin{frame}{El contrato de evento es una API durable}
\UASlideTag{NÚCLEO}{UAIngNavy}
\begin{itemize}
\item Queda retenido y puede ser leído por consumidores desconocidos al publicar.
\item Cambiar semántica rompe replay histórico, no solo clientes actuales.
\item Agregar campos opcionales suele ser más seguro que renombrar o reinterpretar.
\item Debe existir owner, versión, compatibilidad y política de deprecación.
\end{itemize}
\begin{uakeybox}
El desacoplamiento temporal aumenta la responsabilidad sobre el contrato, no la elimina.
\end{uakeybox}
\end{frame}

\UASectionSlide{Sagas: consistencia de negocio\\sin rollback global}

\begin{frame}{Saga de compra: camino feliz}
\UASlideTag{NÚCLEO}{UAIngNavy}
\begin{center}
\begin{tikzpicture}[node distance=0.45cm]
\node[uaBoxO,text width=2.1cm] (a) {1. Reservar\\objeto};
\node[uaBox,text width=2.1cm,right=of a] (b) {2. Autorizar\\pago};
\node[uaBoxO,text width=2.1cm,right=of b] (c) {3. Transferir\\propiedad};
\node[uaBox,text width=2.1cm,right=of c] (d) {4. Notificar\\jugador};
\draw[uaArrow] (a)--(b); \draw[uaArrowO] (b)--(c); \draw[uaArrow] (c)--(d);
\end{tikzpicture}
\end{center}
\begin{block}{Decisión previa}
¿Qué pasos son críticos? ¿Cuáles pueden reintentarse? ¿Qué estado intermedio puede ver el usuario?
\end{block}
\end{frame}

\begin{frame}{Saga de compra: falla y compensación}
\UASlideTag{NÚCLEO}{UAIngNavy}
\begin{center}
\begin{tikzpicture}[node distance=0.45cm]
\node[uaBoxO,text width=2.1cm] (a) {Reserva OK};
\node[uaBox,text width=2.1cm,right=of a] (b) {Pago OK};
\node[uaBoxO,text width=2.1cm,right=of b] (c) {Transferencia\\FALLA};
\node[uaSoft,text width=2.2cm,below=0.9cm of b] (cb) {Reembolso};
\node[uaSoft,text width=2.2cm,below=0.9cm of a] (ca) {Liberar reserva};
\draw[uaArrow] (a)--(b); \draw[uaArrowO] (b)--(c);
\draw[uaArrowO,dashed] (c)--(cb); \draw[uaArrow,dashed] (cb)--(ca);
\end{tikzpicture}
\end{center}
\begin{uawarn}
El reembolso es una transacción nueva. Puede tardar, fallar o dejar movimientos visibles. La saga necesita estados, retries e intervención manual.
\end{uawarn}
\end{frame}

\begin{frame}{Compensación no es rollback}
\UASlideTag{NÚCLEO}{UAIngNavy}
\begin{center}
\begin{tabular}{p{0.27\textwidth}p{0.30\textwidth}p{0.30\textwidth}}
\toprule
Efecto & ¿Se puede ``deshacer''? & Compensación real \\
\midrule
Reserva interna & casi siempre & liberar reserva \\
Cobro & no borrar historial & reembolso \\
Email enviado & imposible retirar & enviar corrección \\
Paquete despachado & efecto físico & devolución / soporte \\
\bottomrule
\end{tabular}
\end{center}
\begin{block}{Lección}
Diseñar una saga exige semántica de negocio, no solo encadenar endpoints.
\end{block}
\end{frame}

\begin{frame}{Coreografía: cada servicio escucha y reacciona}
\UASlideTag{EXTENSIÓN}{UAIngMuted}
\begin{center}
\begin{tikzpicture}[node distance=0.65cm]
\node[uaBoxO] (o) {Orders};
\node[uaBox,right=1.2cm of o] (p) {Payments};
\node[uaBoxO,right=1.2cm of p] (i) {Inventory};
\node[uaBox,right=1.2cm of i] (n) {Notify};
\draw[uaArrow] (o)--node[above,font=\scriptsize]{OrderCreated}(p);
\draw[uaArrowO] (p)--node[above,font=\scriptsize]{PaymentOK}(i);
\draw[uaArrow] (i)--node[above,font=\scriptsize]{ItemMoved}(n);
\end{tikzpicture}
\end{center}
\begin{itemize}
\item Menos coordinador central.
\item Más lógica global distribuida entre contratos.
\item Difícil ver el estado de la saga sin observabilidad adicional.
\end{itemize}
\end{frame}

\begin{frame}{Orquestación: un coordinador conserva el proceso}
\UASlideTag{EXTENSIÓN}{UAIngMuted}
\begin{center}
\begin{tikzpicture}[node distance=0.65cm]
\node[uaBoxO,text width=2.2cm] (orch) {Saga orchestrator\\estado del flujo};
\node[uaBox,above right=0.3cm and 1.3cm of orch] (p) {Payments};
\node[uaBox,right=1.3cm of orch] (i) {Inventory};
\node[uaBox,below right=0.3cm and 1.3cm of orch] (n) {Notify};
\draw[uaArrow] (orch)--(p); \draw[uaArrowO] (orch)--(i); \draw[uaArrow] (orch)--(n);
\end{tikzpicture}
\end{center}
\begin{itemize}
\item Flujo y estados más visibles.
\item Acoplamiento al proceso central.
\item El orquestador debe ser durable, recuperable e idempotente.
\end{itemize}
\end{frame}

\begin{frame}{Elegir coreografía u orquestación}
\UASlideTag{EXTENSIÓN}{UAIngMuted}
\begin{center}
\begin{tabular}{p{0.22\textwidth}p{0.31\textwidth}p{0.31\textwidth}}
\toprule
Criterio & Coreografía & Orquestación \\
\midrule
Flujo & emergente & explícito \\
Visibilidad & requiere reconstrucción & estado central del proceso \\
Acoplamiento & contratos de eventos & al orquestador \\
Complejidad & crece con participantes & crece en coordinador \\
\bottomrule
\end{tabular}
\end{center}
\begin{block}{Regla}
No hay ganador universal. Importan longitud del proceso, necesidad de intervención, auditabilidad y frecuencia de cambio.
\end{block}
\end{frame}

\UASectionSlide{Migrar y operar: una frontera\\debe sobrevivir producción}

\begin{frame}{Strangler Fig: reemplazar por rutas, no por fe}
\UASlideTag{EXTENSIÓN}{UAIngMuted}
\begin{center}
\begin{tikzpicture}[node distance=0.65cm]
\node[uaBoxO,text width=2.3cm] (gw) {Router / gateway};
\node[uaBox,text width=2.6cm,above right=0.45cm and 1.3cm of gw] (old) {Monolito\\rutas restantes};
\node[uaBoxO,text width=2.6cm,below right=0.45cm and 1.3cm of gw] (new) {Nuevo servicio\\ruta extraída};
\draw[uaArrow] (gw)--node[above,font=\scriptsize]{legacy}(old);
\draw[uaArrowO] (gw)--node[below,font=\scriptsize]{migrada}(new);
\end{tikzpicture}
\end{center}
\begin{itemize}
\item Migración incremental y reversible.
\item Permite medir beneficio antes de extraer todo.
\item Requiere resolver sincronización de datos y rutas durante coexistencia.
\end{itemize}
\end{frame}

\begin{frame}{Checklist de readiness antes de distribuir}
\UASlideTag{NÚCLEO}{UAIngNavy}
\begin{columns}[T]
\column{0.48\textwidth}
\begin{itemize}
\item CI/CD y rollback;
\item logs y tracing;
\item métricas y SLO;
\item ownership y on-call;
\item contratos versionados.
\end{itemize}
\column{0.48\textwidth}
\begin{itemize}
\item secretos e identidad;
\item timeouts y retries;
\item pruebas de falla;
\item capacidad de datos;
\item runbooks e incidentes.
\end{itemize}
\end{columns}
\begin{uawarn}
Sin estas capacidades, la separación puede aumentar el tiempo de entrega y reducir confiabilidad aunque el diagrama parezca moderno.
\end{uawarn}
\end{frame}

\begin{frame}{Métricas que demuestran autonomía real}
\UASlideTag{NÚCLEO}{UAIngNavy}
\begin{center}
\begin{tikzpicture}[node distance=0.18cm]
\node[uaBoxO,text width=1.98cm] (a) {Lead time\\por cambio};
\node[uaBox,text width=1.98cm,right=of a] (b) {Deploys\\independientes};
\node[uaBoxO,text width=1.98cm,right=of b] (c) {Blast radius\\incidentes};
\node[uaBox,text width=1.98cm,right=of c] (d) {p99 / SLO\\por servicio};
\node[uaBoxO,text width=1.98cm,right=of d] (e) {Coordinaciones\\entre equipos};
\end{tikzpicture}
\end{center}
\begin{block}{Pregunta}
Si después de extraer todos despliegan juntos y comparten la misma base, ¿qué autonomía se ganó realmente?
\end{block}
\end{frame}

\begin{frame}{Actividad de 15 minutos: diseñar Asteria sin slogans}
\UASlideTag{ACTIVIDAD}{UAIngOrange}
\begin{columns}[T]
\column{0.48\textwidth}
\begin{block}{Elegir una frontera}
\begin{itemize}
\item inventario;
\item combate;
\item ranking;
\item pagos;
\item chat.
\end{itemize}
\end{block}
\column{0.48\textwidth}
\begin{block}{Defenderla}
\begin{itemize}
\item capacidad y owner;
\item invariante local;
\item RPC/evento;
\item falla nueva;
\item autonomía y métrica.
\end{itemize}
\end{block}
\end{columns}
\begin{uainfo}
Una respuesta válida puede mantener módulos en un monolito si demuestra que separar no compra suficiente valor.
\end{uainfo}
\end{frame}

\begin{frame}{Método reusable para decidir una frontera}
\UASlideTag{NÚCLEO}{UAIngNavy}
\begin{center}
\begin{tikzpicture}[node distance=0.35cm]
\node[uaBoxO,text width=2.1cm] (d) {1. Dominio\\e invariante};
\node[uaBox,text width=2.1cm,right=of d] (o) {2. Owner\\y equipo};
\node[uaBoxO,text width=2.1cm,right=of o] (c) {3. Contrato\\y datos};
\node[uaBox,text width=2.1cm,right=of c] (f) {4. Falla y\\operación};
\node[uaBoxO,text width=2.1cm,right=of f] (m) {5. Métrica y\\revisión};
\draw[uaArrow] (d)--(o); \draw[uaArrowO] (o)--(c); \draw[uaArrow] (c)--(f); \draw[uaArrowO] (f)--(m);
\end{tikzpicture}
\end{center}
\begin{block}{Marco del curso}
$D=(P,F,G,S,T)$: la tecnología aparece después del problema, la falla y la garantía.
\end{block}
\end{frame}

\begin{frame}{Volvemos al diagnóstico inicial}
\UASlideTag{ACTIVIDAD}{UAIngOrange}
\small
\begin{enumerate}
\item \textbf{Falso:} un monolito puede ser modular y tener fronteras claras.
\item \textbf{Falso:} un microservicio representa una capacidad y autonomía, no una clase remota.
\item \textbf{Falso:} guardar y publicar son dos escrituras; una falla entre ambas crea inconsistencia.
\item \textbf{Falso:} compensar es ejecutar una nueva acción, no borrar el pasado.
\item \textbf{Falso:} cadenas de dependencias pueden reducir disponibilidad y aumentar latencia.
\end{enumerate}
\end{frame}

\begin{frame}{Síntesis visual: una frontera debe pagar su alquiler}
\UASlideTag{NÚCLEO}{UAIngNavy}
\begin{center}
\begin{tikzpicture}[node distance=0.35cm]
\node[uaBoxO,text width=2.1cm] (a) {Capacidad de\\negocio};
\node[uaBox,text width=2.1cm,right=of a] (b) {Owner e\\invariante};
\node[uaBoxO,text width=2.1cm,right=of b] (c) {Contrato y\\datos};
\node[uaBox,text width=2.1cm,right=of c] (d) {Falla y\\operación};
\node[uaBoxO,text width=2.1cm,right=of d] (e) {Autonomía\\medida};
\draw[uaArrow] (a)--(b); \draw[uaArrowO] (b)--(c); \draw[uaArrow] (c)--(d); \draw[uaArrowO] (d)--(e);
\end{tikzpicture}
\end{center}
\begin{uakeybox}
Distribuir no es progreso por sí mismo. Una frontera vale cuando compra autonomía, escala o aislamiento suficiente para pagar su prima distribuida.
\end{uakeybox}
\end{frame}

\begin{frame}{Exit ticket}
\UASlideTag{ACTIVIDAD}{UAIngOrange}
Responda en cuatro líneas cada punto:
\begin{enumerate}
\item ¿Qué diferencia un monolito modular de una bola de barro?
\item ¿Qué evidencia justifica extraer un microservicio?
\item ¿Cómo resuelve outbox el dual write y qué problema deja abierto?
\item ¿Por qué una compensación de saga no es rollback?
\end{enumerate}
\end{frame}

\begin{frame}{Fuentes para profundizar}
\UASlideTag{NÚCLEO}{UAIngNavy}
\begin{itemize}
\item Kleppmann, \emph{Designing Data-Intensive Applications}: servicios, eventos, consistencia y stream processing.
\item Martin Fowler, \emph{Monolith First}, \emph{Microservice Trade-Offs} y \emph{Strangler Fig}.
\item Shopify Engineering: modular monolith, componentización y Packwerk.
\item AWS Prescriptive Guidance: transactional outbox, saga choreography/orchestration y strangler fig.
\item Netflix TechBlog: casos de microservicios, pipelines y event-driven architecture.
\end{itemize}
\end{frame}

\end{document}
